\section{Uniform arithmetic saturation of the dyadic word image}
\label{sec:fr-uniform-dyadic-lie}

The characteristic-two generators admit a uniform arithmetic
construction. The useful multiplier is $-3$ in the $2$-adic
component: its power substitution has no quadratic or cubic
correction modulo two.

Put $M=\ker\lambda$, $H_M=H\cap M$, and
\[
 F^nM=\{g\in M:f_g-1\text{ has no dyadic Magnus term of degree }<n\}.
\]
Every element of $M$ belongs to $F^3M$. Let
\[
 R=\bigcap_{n\ge3}F^nM
   =\ker\bigl(M\longrightarrow
                \operatorname{Aut}(\widehat F_2^{(2)})\bigr),
 \qquad R_D=D\cap R.
\]
The equality in this definition is justified below.

\begin{theorem}\label{thm:fr-uniform-dyadic-saturation}
For every degree $n\ge3$, the full and arithmetic leading images
are the entire integral Grothendieck--Teichm\"uller Lie lattice:
\[
 \operatorname{in}_n(F^nM)
 =\operatorname{in}_n(F^nM\cap H)
 =\mathcal L_n,
 \qquad
 \mathcal L_n=
 \mathfrak{grt}_1(\mathbf Q_2)_n
          \cap\operatorname{Lie}_{\mathbf Z_2}(S,T)_n.
\]
There are arithmetic elements $z_{2r+1}\in H_M\cap F^{2r+1}M$
whose reduced leading terms are precisely the characteristic-two
free generators of Willwacher's Theorem~1.3. Their leading terms
freely generate $\mathcal L=\bigoplus_n\mathcal L_n$ over
$\mathbf Z_2$.

One has
\[
 M=RH_M,\qquad G=RH,\qquad G=R_DH.
\]
In particular the full-GT and arithmetic images in
$\operatorname{Aut}(\widehat F_2^{(2)})$ coincide.
The normalized subgroup $R_D$ is closed and normal, and
\[
 B\simeq R_D/\overline{[G,R_D](R_D\cap H)}.
\]
The kernel of restriction from $X_2=H^1(N,\mathbf F_2)^G$ to
$R_D$ is exactly
\[
 \left(\bigoplus_{q\ {\rm odd\ prime}}\mathbf F_2\psi_q\right)
                   \oplus\mathbf F_2\eta_3.
\]
Moreover $R_D$ is a proper subgroup of infinite index in
$E_{15}=C_\infty\cap F^{15}M$.
\end{theorem}

The proof uses the defining-equation and characteristic-two
freeness results of
\cite[equations (1)--(6), p.~1; \S1.1 and Theorem~1.3, p.~2;
\S4.3, p.~9]{Willwacher}. It does not require an a priori
arithmetic realization of a rational or integral Lie solution.

\begin{lemma}\label{lem:fr-integral-leading-basis}
The finite-rank module $\mathcal L_n$ is saturated in
$\operatorname{Lie}_{\mathbf Z_2}(S,T)_n$, and reduction gives
an injection
\[
 \mathcal L_n/2\mathcal L_n
       \hookrightarrow\mathfrak{grt}_1(\mathbf F_2)_n.
\]
If a collection of elements of $\mathcal L_n$ reduces to a
basis of the space on the right, it is a $\mathbf Z_2$-basis
of $\mathcal L_n$.
\end{lemma}

\begin{proof}
The intersection with a rational subspace is saturated:
if $2v$ belongs to it and $v$ is integral, then $v$ belongs
to the same rational subspace. The integral defining equations,
including cabling, remain valid after reduction modulo two.
Their finite integral matrix description is recalled in
\cite[\S1.1, p.~2]{Willwacher}. Saturation makes the displayed
reduction map injective.

For the final assertion, the proposed reductions are linearly
independent in $\mathcal L_n/2\mathcal L_n$, and their images
span its ambient target. They therefore form a basis of
$\mathcal L_n/2\mathcal L_n$. Nakayama's lemma makes their lifts
generate $\mathcal L_n$; their independence follows by
successive reduction modulo powers of two, or by equality of
the finite ranks.
\end{proof}

The necessary leading-map lemma applies here in every degree.
For completeness, project the full profinite pentagon to the
pro-$2$ and then rational prounipotent completions, as in
\cite[\S3, p.~554, immediately before Corollary~12]{Furusho2010}.
The substitution $x=e^S$, $y=e^T$ preserves the first nonzero
term. That term is primitive and hence is a free Lie polynomial
over $\mathbf Q_2$. Taking the first nonzero degree in the group
pentagon gives the linear pentagon in
$\mathfrak t_4(\mathbf Q_2)$: the degree-one class of a pure
braid generator is $t_{ij}$, products add these classes, and
products of two nonconstant word terms have degree at least
$2n>n$. Duality and the hexagon likewise give antisymmetry
and the three-cycle identity. Cabling is redundant in
characteristic zero
\cite[p.~1, after equation~(6)]{Willwacher}.
Thus the leading term belongs to
$\mathfrak{grt}_1(\mathbf Q_2)_n$.
Its integral associative coefficients are integral Lyndon
coefficients, because the Lyndon expansions are unit triangular.
It consequently belongs to $\mathcal L_n$. This is only a
necessary inclusion; the following construction supplies the
arithmetic realization.

\begin{lemma}\label{lem:fr-minus-three-jet}
Choose $h\in H$ with $t=\lambda(h)$ and $t_2=-3$.
Put $C=[S,T]$. For $z\in F^nM$, $n\ge3$, write
$q=\operatorname{in}_n(z)$ and define, using the full
profinite exponent,
\[
 \Delta_n(z)=hzh^{-1}z^{-t^n}.
\]
Then $\Delta_n(z)\in F^{n+1}M$, and its Magnus series satisfies
\begin{equation}\label{eq:fr-minus-three-jet}
 f_{\Delta_n(z)}
 \equiv1+\{C,\overline q\}
        \pmod{\,2,\ (S,T)^{n+3}}.
\end{equation}
In particular every coefficient through degree $n+1$ is
even, while its degree-$n+2$ term modulo two is the indicated
Ihara bracket.
\end{lemma}

\begin{proof}
We specify the ambient convention to check the jet precisely.
For a pair $(s,f)$ let
\[
 \Phi_{s,f}(x)=x^s,\qquad
 \Phi_{s,f}(y)=f^{-1}y^sf.
\]
Composition gives
\[
 (s,f)(u,g)=(su,\ f\,\Phi_{s,f}(g)).
\]
Consequently, in the ambient group of these automorphisms,
\[
 h=(1,f_h)(t,1).
\]
The power pair $(t,1)$ is an automorphism, with inverse
$(t^{-1},1)$, so this factorization is legitimate even though
its factors need not individually belong to GT.

The degree-two coefficient of $f_h$ is $\tfrac13 C$.
Indeed, put $m=(t_2-1)/2=-2$ in the multiplier hexagon
and compute its logarithm through degree two. The three
$f_h$ terms contribute $3c_2(h)C$, and the intervening
product $z^my^mx^m$ contributes
$-\tfrac12m(m+1)C$. Thus
$c_2(h)=m(m+1)/6=1/3$. In particular its reduction is $C$.

Conjugation by $(t,1)$ replaces the variables of $f_z$ by
$x^t,y^t$. Modulo two,
\[
 (1+S)^{-3}-1=S+O(S^4),\qquad
 (1+T)^{-3}-1=T+O(T^4),
\]
since $\binom{-3}{2}=6$ and $\binom{-3}{3}=-10$ are even.
Therefore this substitution leaves every term of $f_z$
through degree $n+2$ unchanged modulo two, including its
unknown terms in degrees $n+1$ and $n+2$.

For multiplier-one ambient pairs with leading terms $A$ and
$Q$ of degrees $d$ and $n$, the first term of their group
commutator is
\[
 \{A,Q\}=[A,Q]+\partial_AQ-\partial_QA,\qquad
 \partial_A(S)=0,\quad\partial_A(T)=[T,A].
\]
This follows directly by subtracting the two products
$a\,\Phi_a(b)$ and $b\,\Phi_b(a)$: in degree $d+n$ their
difference is $AQ-QA+\partial_AQ-\partial_QA$.
Taking $d=2$ shows that conjugation by $(1,f_h)$ adds exactly
$\{C,\overline q\}$ in degree $n+2$ modulo two. Higher terms
of $f_h$ contribute only in degrees at least $n+3$.

Finally, through degree $n+2$, a profinite power of $z$
has series $1+r_2(f_z-1)$: all product and Ihara
self-substitution corrections begin in degree $2n>n+2$.
Since $t_2^n$ is odd, multiplication by $z^{-t^n}$ cancels
the unchanged $f_z-1$ jet modulo two. This proves
\eqref{eq:fr-minus-three-jet}. Integrally, the degree-$n$
term cancels as well, because conjugation has weight
$t_2^n$ on the leading quotient; hence
$\Delta_n(z)\in F^{n+1}M$.
\end{proof}

\begin{proposition}\label{prop:fr-all-odd-arithmetic-generators}
There exist elements $z_{2r+1}\in H_M\cap F^{2r+1}M$
whose leading reductions are the generators
\[
 s_3=[S+T,[S,T]],\qquad
 s_{2r+3}=\{C,s_{2r+1}\}
\]
of $\mathfrak{grt}_1(\mathbf F_2)$.
In each degree the leading terms of the corresponding
arithmetic Hall commutator words form a
$\mathbf Z_2$-basis of $\mathcal L_n$.
\end{proposition}

\begin{proof}
Take the established arithmetic $\sigma\in H_M$ with
$a=m_3(\sigma)$ odd, and put $z_3=\sigma$.
Its leading reduction is $s_3$. Theorem~1.3 of
\cite[p.~2]{Willwacher} says that the $s_{2r+1}$ freely
generate the characteristic-two Lie algebra under the
Ihara bracket.

Suppose arithmetic generators have been constructed through
the odd degree $n$, and their leading reductions are the
specified generators. Every free Hall word of even weight
$n+1$ is a bracket word in lower odd generators; no new
generator has that even weight. Realize these words by
group commutators in the already constructed $z_j$.
Their leading reductions form the full free Hall basis of
$\mathfrak{grt}_1(\mathbf F_2)_{n+1}$.
Lemma~\ref{lem:fr-integral-leading-basis} makes their leading
terms a $\mathbf Z_2$-basis of $\mathcal L_{n+1}$.
When this characteristic-two space is zero, the integral
space is zero as well.

By Lemma~\ref{lem:fr-minus-three-jet},
$\Delta_n(z_n)$ starts in degree $n+1$ and its degree-$n+1$
term is even. Saturation implies that this term lies in
$2\mathcal L_{n+1}$. Express it in the arithmetic Hall basis
just constructed. All coordinates are even.
Multiply $\Delta_n(z_n)$ by inverse profinite powers of
those Hall words, with these even $2$-adic exponents, to
kill its whole degree-$n+1$ term.
Let the result be $z_{n+2}$.

Each correcting word starts in degree $n+1$. Its even
power is trivial modulo two through degree $n+2$, because
$2(n+1)>n+2$. Products of the corrections have the same
property. Thus the correction does not change the
degree-$n+2$ term in
\eqref{eq:fr-minus-three-jet}. The new arithmetic word
starts in degree $n+2$ and has leading reduction
$\{C,s_n\}=s_{n+2}$.

In odd degree $n+2$, this new generator together with all
Hall bracket words in lower generators reduces to the
complete characteristic-two Hall basis. Apply
Lemma~\ref{lem:fr-integral-leading-basis} once more.
This proves the induction, both in even and odd degrees.
It also proves that the free Lie algebra over $\mathbf Z_2$
on the arithmetic leading terms of the $z_{2r+1}$ maps
isomorphically to $\mathcal L$, degree by degree.
\end{proof}

We record two properties of this construction needed for the
normalized conclusion. Every correcting Hall word in the
induction is a commutator in $M$. Therefore all its rational
Kummer coordinates vanish and its dyadic beta multiplier is
one. If $c_3=1$ and
\[
 c_{n+2}=(t-t^n)c_n\in\widehat{\mathbf Z},
\]
then, for every rational prime including two,
\begin{equation}\label{eq:fr-uniform-kummer-collinearity}
 \rho_q(z_n)=c_n\rho_q(\sigma).
\end{equation}
The same beta-trivial corrections leave every beta moment
unchanged. Hence $z_{2r+1}$ has zero moments in preceding
odd weights, and its depth-one coefficient is
\begin{equation}\label{eq:fr-uniform-primitive-depth}
 [S^{2r}T]f_{z_{2r+1}}
 =\frac{2\prod_{j=1}^{r-1}
       \bigl((-3)^{2r+1}-(-3)^{2j+1}\bigr)}{(2r)!}
                         \tau_{2r+1}(\sigma)\ne0.
\end{equation}
The established all-$M$ moment parity makes the last moment
odd. In particular the depth-one coefficient is nonzero on
each odd integral generator. Its valuation is
\[
 3(r-1)+v_2((r-1)!)+1-v_2((2r)!)
                   =2r-2-v_2(r).
\]
The empty product for $r=1$ has the same interpretation.

\begin{proof}[Proof of Theorem~\ref{thm:fr-uniform-dyadic-saturation}]
Every full leading term belongs to $\mathcal L_n$.
Proposition~\ref{prop:fr-all-odd-arithmetic-generators}
realizes an entire integral basis there by arithmetic
words. Their profinite powers and products realize all
integral linear combinations. This proves the two leading
image equalities.

For any $g\in M$, correct its leading terms successively
by these arithmetic basis words. For each $n$ this gives
an $h_n\in H_M$ with $g h_n^{-1}\in F^nM$.
The sets
\[
 \{h\in H_M:g h^{-1}\in F^nM\}
\]
are nonempty, closed, and nested in the compact group
$H_M$. Their intersection is nonempty. Thus $M=RH_M$.
The cyclotomic multiplier of $H$ is all of
$\widehat{\mathbf Z}^{\times}$, so $G=MH=RH$.

The pro-$2$ Magnus embedding is faithful. Indeed, for every
finite $2$-group quotient $Q$, the augmentation ideal of
$\mathbf F_2[Q]$ is nilpotent. This follows by induction using
a central element $c$ of order two: the ideal generated by $c-1$
has square zero, and quotienting by it gives the group algebra
of $Q/\langle c\rangle$. Substitution of $x-1,y-1$ in this
finite group algebra therefore factors through a finite Magnus
truncation. A word with trivial Magnus series is consequently
trivial in every finite $2$-group quotient, hence in the free
pro-$2$ group. The intersection of the word filtration is thus
exactly the condition that the pro-$2$ word $f_g$ be one. This is also the
kernel of the multiplier-one action on the free pro-$2$
group. Indeed, if that action is the identity, $f_g$
commutes with $y$. In the Magnus algebra the centralizer
of $T$ consists of series in $T$ alone; and $f_g(1,y)=1$
because $f_g$ is in the commutator subgroup. Therefore
$f_g=1$. The converse is immediate. This justifies the
definition of $R$ and the equality of full and arithmetic
dyadic action images. It also shows that $R$ is closed
and normal in $G$.

We now preserve the full rational Kummer normalization.
Let $U$ be the closed subgroup of $H_M$ generated by all
$z_{2r+1}$, and put
\[
 A=\overline{[U,U]}.
\]
Both beta and every rational Kummer coordinate are
homomorphisms on $M$, so $A\subset H\cap C_\infty$.
Moreover \eqref{eq:fr-uniform-kummer-collinearity} gives
$z_n\sigma^{-c_n}\in D$.
In the quotient by $\overline{[G,D]}$, the image of $D$
is central, and all the $z_n$ are central factors times
profinite powers of the same $\sigma$. Their images
therefore commute. By continuity,
\begin{equation}\label{eq:fr-arithmetic-derived-in-relative-commutator}
 A\subset\overline{[G,D]}\subset\overline{[G,N]}.
\end{equation}

In the free integral Lie algebra on the leading
$z_{2r+1}$, the bracket submodule has a Hall basis
consisting of words of generator length at least two.
Every such bracket has zero depth-one coefficient.
In an odd degree the remaining abelianization is
rank one, and its generator has nonzero depth-one
coefficient by
\eqref{eq:fr-uniform-primitive-depth}.
It follows, integrally, that the kernel of the
depth-one functional on $\mathcal L_n$ is exactly
the bracket submodule; in even degrees the whole
space is that submodule.

For an element of $C_\infty$, all depth-one coefficients
vanish because its dyadic beta multiplier is one.
Its leading term can consequently be corrected entirely
by arithmetic Hall bracket words, all lying in $A$.
Repeating this correction degree by degree remains
inside $C_\infty$. For each $c\in C_\infty$, the closed
sets
\[
 \{a\in A:c a^{-1}\in F^nM\}
\]
are nonempty and nested in the compact group $A$.
Their intersection gives
\[
 C_\infty=(C_\infty\cap R)A=R_DA.
\]
The last equality uses the fact that a trivial dyadic
word has trivial dyadic beta multiplier, so
$C_\infty\cap R=D\cap R$.
The established factorization $G=C_\infty H$ now gives
$G=R_DH$. To check the displayed presentation of $B$ explicitly,
put $K=N\cap H$ and $C_R=\overline{[G,R_D]}$.
Then $N=R_DK$. In $G/C_R$, the image of $R_D$ is central
and $G=R_DH$, so the image of $[G,N]$ is contained in the
image of $K$. Hence
\[
 \overline{[G,N]K}=C_RK.
\]
These products are closed because $C_R$ is closed and normal
and $K$ is compact. Since $C_R\subset R_D$,
\[
 R_D\cap C_RK=C_R(R_D\cap H).
\]
Restriction of the quotient map from $N$ therefore yields
the claimed canonical isomorphism for $B$.

Every $\chi\in X_2$ kills $A$ by
\eqref{eq:fr-arithmetic-derived-in-relative-commutator}.
Thus $\chi$ vanishes on $R_D$ if and only if it vanishes
on $C_\infty$. The already computed restriction kernel
for $C_\infty$ is precisely the rational Kummer span
plus $\mathbf F_2\eta_3$, proving the asserted kernel.

Finally, the arithmetic word
$[z_3,[z_3,z_9]]$ has a nonzero primitive leading term
of degree fifteen, since its reduction is a nonzero
free Hall word. It lies in $A\cap F^{15}M\subset E_{15}$.
Its procyclic image modulo $R_D$ is infinite: multiplying
its nonzero leading coefficient by a nonzero $2$-adic
exponent cannot give zero. Hence $R_D$ is proper of
infinite index in $E_{15}$.
\end{proof}

\begin{remark}
The conclusion concerns the dyadic free-group action and
the normalized domain carrying the relative obstruction.
It does not identify $R_D$ with the trivial group or
prove that $B$ has no nonzero quadratic characters.
\end{remark}
